\documentclass[11pt]{article}
\usepackage[margin=1.1in]{geometry}
\usepackage{times}
\usepackage[T1]{fontenc}
\usepackage{setspace}
\onehalfspacing
\usepackage{parskip}
\pagenumbering{gobble}

\title{The Cost of Not Building It}
\author{Flyxion}
\date{September 2026}

\begin{document}
\maketitle

Every objection to moving faster is met with an accounting exercise, and the exercise is rigged before the first number is entered, because the two columns are not held to the same standard of evidence.

In one column sit the benefits: diseases cured, poverty ended, abundance distributed to everyone who currently lacks it, all attributed to a system that does not yet exist, built to a specification that has not been finalized, on a timeline that has already been wrong before. These benefits are not discounted for uncertainty. They are counted at full value, presented as lives that will be lost if the timeline slips, which quietly converts a speculative upside into a moral debt owed by anyone urging caution, as though delay were the active harm and the undeployed system the neutral default.

In the other column sit the costs: the errors the current system is already producing, the jobs already displaced, the harms already documented in the world as it presently exists. These are treated as temporary. Implementation details. Bugs on a road that leads somewhere better, mentioned in the passive voice and scheduled for a fix in the next version, their weight discounted precisely because they are real and therefore bounded, unlike the imagined benefits, which are unbounded because nothing constrains a number that describes a world that has not happened yet.

This is not a disagreement about how much weight the future deserves relative to the present, which would be a genuine and difficult question. It's a switch in method depending on which direction the number points. Future benefit gets counted as certain. Present harm gets counted as provisional. Apply that asymmetry to any other domain and the absurdity is immediate: a pharmaceutical company arguing that a drug's currently observed side effects are minor implementation details, while its projected future benefit to global health, once the formulation is perfected, should be treated as an already-banked certainty, would not be granted the accounting. It would be asked why the same evidentiary standard doesn't apply to both entries in the ledger, and it would not have a good answer, because the answer is that one standard makes the argument work and the other doesn't.

\end{document}
